% !PRINTED 426
\weberpage{426}

%<<CONT
can be made, the fundamentals of which may fittingly be inserted
here, where the analogy with the theory of the roots of unity
shows itself clearly.

So let $n$ now be a \gesperrt{prime} and
\weq{1}{f(x) = a_0 x^m + a_1 x^{m-1} + \cdots + a_{m-1} x + a_m}
an integral function of $x$, whose coefficients $a$ are integers,
and $a_0$ not divisible by $n$. If we put for $x$ an integer
$\alpha$ such that $f(\alpha)$ becomes divisible by $n$, then we
say (by analogy with equations) that $\alpha$ is a root of the
congruence of degree $m$:
\weq{2}{f(x) \equiv 0 \quad (\mathrm{mod}\ n).}
If $\alpha$ is increased by a multiple of $n$, it remains a root
of the congruence (2). Such roots, congruent with respect to the
modulus $n$, are not counted as distinct. Under this assumption
we can state the theorem:

\gesperrt{1. A congruence of degree $m$ for a prime
modulus cannot have more than $m$ distinct roots.}

The theorem is true for $m = 1$, for the congruence $a_0 x + a
\equiv 0\ (\mathrm{mod}\ n)$ has only one root, namely, if
% [sic] The source prints modulus m here although the inverse is being taken modulo n; preserved diplomatically.
$a_0'$ is determined so that $a_0 a_0' \equiv 1\ (\mathrm{mod}\ m)$
holds, $\alpha \equiv -a a_0'\ (\mathrm{mod}\ n)$. We therefore
now assume our theorem to be proved for the degree $m-1$ and
deduce its validity for the degree $m$ from this. By \S.\,4 we
can, if at first $\alpha$ is arbitrary, set
\weq{3}{\frac{f(x) - f(\alpha)}{x - \alpha} = f_1(x),}
in which $f_1(x)$ is an integral function of degree $(m-1)$,
which, if $\alpha$ is an integer, has integral
coefficients. If therefore $f(\alpha) \equiv 0$, it follows
from (3) that
\weq{4}{f(x) \equiv (x - \alpha)\, f_1(x) \quad (\mathrm{mod}\ n).}
Every root $\beta$ of the congruence (2) must therefore satisfy
the condition $(\beta - \alpha)\, f_1(\beta) \equiv 0\ (\mathrm{mod}\ n)$.
Hence either $\beta - \alpha$ or $f_1(\beta)$ is
divisible by $n$.

Now by hypothesis there are at most $m-1$ values $\beta$
for which $f_1(\beta)$ becomes divisible by $n$; if therefore
there is still an $m$th root of (4), it must be equal to $\alpha$.

If therefore a congruence of the form $f(x) \equiv 0\ (\mathrm{mod}\
n)$ has more roots than its degree amounts to, then we conclude
that
%CONT>>
